% ============================================================================
%  SECCIÓN 4.11: EVIDENCIAS DEL AVANCE
% ============================================================================

\section{Evidencias del avance}

A continuaci\'on se presentan las evidencias del avance correspondiente al
PDF~4, organizadas seg\'un la gu\'ia: formularios funcionando, validaciones,
mensajes de error y \'exito, cambio de estado, funcionalidad principal,
listas din\'amicas, b\'usqueda y repositorio actualizado.

% --- EVIDENCIA 1 ---
\subsubsection*{Evidencia 1: Formulario funcionando}

El formulario de registro captura los datos introducidos por el usuario.
Los valores ingresados se conservan en el estado del formulario mientras
se completa:

\begin{figure}[H]
  \centering
  \caption{Evidencia: Formulario de registro en funcionamiento}
  \label{fig:evidencia4-formulario}
  \contenidoConFuente{%
    \includegraphics[width=0.28\linewidth]{evidencia-formulario-registro.jpg}%
  }{Elaboraci\'on propia (captura de la aplicaci\'on).}
\end{figure}

% --- EVIDENCIA 2 ---
\subsubsection*{Evidencia 2: Validaci\'on de campos}

Al enviar el formulario con datos incompletos o inv\'alidos, cada campo
incumplido muestra su mensaje de error espec\'ifico debajo, con borde rojo:

\begin{figure}[H]
  \centering
  \caption{Evidencia: Validaci\'on de campos del registro}
  \label{fig:evidencia4-validacion-registro}
  \contenidoConFuente{%
    \includegraphics[width=0.28\linewidth]{evidencia-validacion-registro.jpg}%
  }{Elaboraci\'on propia (captura de la aplicaci\'on).}
\end{figure}

% --- EVIDENCIA 3 ---
\subsubsection*{Evidencia 3: Mensaje de error}

En el inicio de sesi\'on, un email sin formato v\'alido es bloqueado por la
validaci\'on antes de llegar al servidor; las credenciales incorrectas se
comunican mediante la caja de error del servidor:

\begin{figure}[H]
  \centering
  \caption{Evidencia: Mensaje de error en el inicio de sesi\'on}
  \label{fig:evidencia4-error-login}
  \contenidoConFuente{%
    \includegraphics[width=0.28\linewidth]{evidencia-validacion-campo-login.jpg}%
  }{Elaboraci\'on propia (captura de la aplicaci\'on).}
\end{figure}

% --- EVIDENCIA 4 ---
\subsubsection*{Evidencia 4: Mensaje de operaci\'on exitosa}

Tras completar el registro con datos correctos, la aplicaci\'on informa que
la cuenta fue creada e indica el siguiente paso (verificar el correo):

\begin{figure}[H]
  \centering
  \caption{Evidencia: Mensaje de operaci\'on exitosa}
  \label{fig:evidencia4-exito-registro}
  \contenidoConFuente{%
    
\includegraphics[width=0.28\linewidth]{evidencia-exito-registro.jpg}%
  }{Elaboraci\'on propia (captura de la aplicaci\'on).}
\end{figure}

% --- EVIDENCIA 5 ---
\subsubsection*{Evidencia 5: Cambio de estado}

El bot\'on de favorito refleja el estado interno: coraz\'on vac\'io cuando el
escenario no est\'a guardado y coraz\'on lleno despu\'es de presionarlo. La
interfaz cambia sin recargar la pantalla:

\begin{figure}[H]
  \centering
  \caption{Evidencia: Cambio de estado del bot\'on de favorito}
  \label{fig:evidencia4-favorito-antes}
  \contenidoConFuente{%
    \includegraphics[width=0.28\linewidth]{evidencia-favorito-antes.jpg}%
  }{Elaboraci\'on propia (captura de la aplicaci\'on).}
\end{figure}

\begin{figure}[H]
  \centering
  \caption{Evidencia: Estado actualizado tras presionar el bot\'on}
  \label{fig:evidencia4-favorito-despues}
  \contenidoConFuente{%
    \includegraphics[width=0.28\linewidth]{evidencia-favorito-despues.jpg}%
  }{Elaboraci\'on propia (captura de la aplicaci\'on).}
\end{figure}

% --- EVIDENCIA 6 ---
\subsubsection*{Evidencia 6: Funcionalidad principal interactiva}

La funcionalidad principal del avance es la gesti\'on de favoritos: el
usuario guarda escenarios deportivos desde la pantalla de detalle y la
lista se sincroniza en tiempo real con la base de datos (evidencias 5 y
figura~\ref{fig:flujo-favoritos}). Esta funcionalidad se relaciona
directamente con el MVP definido en el PDF~1, cuyo prop\'osito es facilitar
el acceso a informaci\'on confiable sobre la infraestructura deportiva.

% --- EVIDENCIA 7 ---
\subsubsection*{Evidencia 7: Lista din\'amica}

La lista de escenarios responde a los cambios de estado: al aplicar
filtros o modificar favoritos, la lista se re-renderiza autom\'aticamente
reflejando la nueva situaci\'on de los datos:

\begin{figure}[H]
  \centering
  \caption{Evidencia: Lista din\'amica de escenarios}
  \label{fig:evidencia4-lista-dinamica}
  \contenidoConFuente{%
    \includegraphics[width=0.28\linewidth]{evidencia-lista-dinamica.jpg}%
  }{Elaboraci\'on propia (captura de la aplicaci\'on).}
\end{figure}

% --- EVIDENCIA 8 ---
\subsubsection*{Evidencia 8: B\'usqueda}

La b\'usqueda filtra los escenarios mientras el usuario escribe, con un
retardo de 300~ms para evitar filtrados intermedios:

\begin{figure}[H]
  \centering
  \caption{Evidencia: B\'usqueda de escenarios por nombre}
  \label{fig:evidencia4-busqueda}
  \contenidoConFuente{%
    \includegraphics[width=0.28\linewidth]{evidencia-busqueda.jpg}%
  }{Elaboraci\'on propia (captura de la aplicaci\'on).}
\end{figure}

Adem\'as, la b\'usqueda puede combinarse con los filtros por tipo de
escenario y por deporte:

\begin{figure}[H]
  \centering
  \caption{Evidencia: B\'usqueda combinada con filtros}
  \label{fig:evidencia4-busqueda-filtros}
  \contenidoConFuente{%
    \includegraphics[width=0.28\linewidth]{evidencia-busqueda-filtros.jpg}%
  }{Elaboraci\'on propia (captura de la aplicaci\'on).}
\end{figure}

% --- EVIDENCIA 9 ---
\subsubsection*{Evidencia 9: Flujo completo de la funcionalidad}

El flujo completo de la funcionalidad de favoritos queda documentado por
las evidencias anteriores: b\'usqueda del escenario
(figura~\ref{fig:evidencia4-busqueda}), apertura del detalle, presi\'on del
bot\'on de guardado (figuras~\ref{fig:evidencia4-favorito-antes}
y~\ref{fig:evidencia4-favorito-despues}) y verificaci\'on del elemento en la
lista de favoritos (figura~\ref{fig:evidencia4-lista-dinamica}). El
diagrama del proceso completo se presenta en la
figura~\ref{fig:flujo-favoritos} de la secci\'on de operaciones CRUD.

% --- EVIDENCIA 10 ---
\subsubsection*{Evidencia 10: Proyecto actualizado en GitHub}

El historial de commits documenta la evoluci\'on progresiva del proyecto,
incluyendo la integraci\'on de l\'ogica, validaciones y estado del avance:

\begin{lstlisting}[style=terminal]
feat(upload): agregar funcionalidad de subir imagenes desde la app
fix(types): agregar events al tipo ScenarioWithDetails para detalle
chore(deps): agregar react-dom y react-native-web para soporte web
docs: agregar metodologia Scrum, historias de usuario y correcciones
docs: agregar contenido del PDF 4 - logica, formularios, validaciones
\end{lstlisting}

La captura general del repositorio se present\'o en la
figura~\ref{fig:repositorio-github} de la configuraci\'on de Git y GitHub.
